\subsection{Supervised Fine Tuning}
\label{subsec:SFT}

% \todo{Wenliang: New data that was added for SFT. Other SFT details}

% To train the student model, starting from the pre-trained base model, we first conduct Supervised Fine Tuning (SFT) using packed sequences of length $294{,}912$ with a batch size of 64. We train for $166{,}400$ samples using a cosine learning-rate scheduler with peak learning rate $1.5e\!-\!5$ and minimum learning rate $1e\!-\!6$, and 9{,}600 warmup samples. As in pre-training, we retain the shared-weight MTP objective during SFT, using two MTP layers with per-token auxiliary loss scaling factor 0.1. 

% To train the student model, we start from the pre-trained base model and perform Supervised Fine-Tuning (SFT) on packed sequences of length $294{,}912$ with a global batch size of $64$. We train on $204{,}800$ samples using a cosine learning-rate schedule with a peak learning rate of $1.5 \times 10^{-5}$, a minimum learning rate of $1 \times 10^{-6}$, and $9{,}600$ warmup samples. 

% % We optimize a sequence-averaged token-level cross-entropy loss. Let the global batch consist of $B$ packed sequences, where sequence $i$ contains $T_i$ supervised (response) tokens with per-token loss $\ell_{i,t}$; prompt tokens are masked out. We first average the loss over the supervised tokens within each sequence, then average uniformly across the $B$ sequences in the batch:
% % \begin{equation*}
% % \mathcal{L} = \frac{1}{B} \sum_{i=1}^{B} \frac{1}{T_i} \sum_{t=1}^{T_i} \ell_{i,t},
% % \end{equation*}
% % Since packs are filled to a nearly constant total length but only response tokens are supervised, this normalization prevents response-dense packs (e.g., short-prompt/long-response data) from dominating the update, yielding a more balanced objective across heterogeneous input/output length ratios.

% As in pre-training, we retain the shared-weight MTP objective during SFT, using two MTP layers with a per-token auxiliary-loss scaling factor of $0.1$. 

To train the student model, we start from the pre-trained base model and perform Supervised Fine-Tuning (SFT) in two stages, following \supermodel~\citep{nvidia2026nemotron3superopen}. In Stage~1, we train on packed sequences of length $294{,}912$ tokens with global batch size $64$ for $204{,}800$ samples, using a cosine learning-rate schedule with peak learning rate $1.5 \times 10^{-5}$, minimum learning rate $1 \times 10^{-6}$, and $9{,}600$ warmup samples. In Stage~2, we extend packed sequence to $515{,}000$ tokens, augmenting the mixture with additional long-context data up to $512$K tokens. We train with global batch size $64$ for $19{,}200$ samples, using the same learning-rate schedule with peak $1 \times 10^{-5}$ to minimum $2 \times 10^{-6}$ and $6{,}400$ warmup samples. 
% Since long-context examples typically contain many prompt/context tokens but relatively few supervised output tokens, Stage~2 uses per-sequence normalization to prevent response-dense shorter examples from dominating the update. This combination of longer packed-sequence training and sample-level normalization better balances short-context reasoning with long-context capability. 
As in pre-training, we retain the shared-weight MTP objective during SFT, using two MTP layers with a per-token auxiliary-loss scaling factor of $0.1$.

\subsubsection{Data}
\paragraph{Long Context.}
We prepare 512K long-context SFT data based on the synthetic data pipeline following~\citep{nvidia2026nemotron3superopen}. 
Our data is aimed at improving long-context abilities including, but not limited to multi-document reasoning, sequential scanning, and querying synthetic tables.
%Our data is aimed at improving the abilities including long-context multi-document reasoning, sequential scanning, and SQL queries. 
\paragraph{Efficiency and Control.}
The SFT data include two components for reasoning efficiency and control. The first is training samples
generated by GPT-OSS-120B in its medium-effort mode on prompts of math reasoning, STEM question answering and instruction following. These SFT data initiate Ultra's medium-effort mode which is later optimized during the RLVR stage. The second component is training samples where the reasoning traces are truncated to random reasoning budgets while the responses remain the same. This is similar to Nemotron 3 Nano and Super with one design change: the \texttt{</think>} tokens in the truncated samples are masked from SFT training loss.


\paragraph{Safety.}
To instill robust safety behaviors, first, we retain the $45$K safety data blend curated in \supermodel~\citep{nvidia2026nemotron3superopen} with prompts curated from diverse sources and responses synthetically generated conditioned on a response policy mapped to each prompt. Our two-stage response and reasoning generation framework enables deliberate reflection on safety guidelines in the reasoning traces and ensures that the final responses are consistent, policy-compliant, and contextually appropriate. Further, to curate a multilingual set for \ourmodel training, we translated the \supermodel safety data blend into six languages — German, Spanish, French, Japanese, Italian, and Chinese - using NeMo Skills' chunked translation capability, ensuring sentence-by-sentence parity, and paired with NVIDIA Riva Translate 4B v1.1 as the translation backbone. 

To improve translation quality, each translated example was back-translated into English and compared against the original English prompt-response pair. Examples with semantic similarity below 0.8 were filtered out, removing approximately $\approx10 - 15\%$ of examples per translated language. The highest- and lowest-scoring translations were manually spot-checked to verify that this filter was removing examples with clear translation failures or structural issues. After filtering, we used stratified sampling to create a balanced dataset across the six translated languages. The final safety blend contains about $\approx135$K samples: $\approx45$K in English and $\approx15$K from each of the translated languages.
% in \texttt{de}, \texttt{es}, \texttt{fr}, \texttt{ja}, \texttt{it}, and \texttt{zh}.

\paragraph{Search Capabilities.}
During SFT, we expose the model to search data spanning different difficulty levels and tool-use harnesses to improve generalization.

First, we retain the search trajectories from the \supermodel~\citep{nvidia2026nemotron3superopen} dataset, whose seed prompts are grounded in the Wikidata knowledge graph. Specifically, the dataset selects well-connected hub entities and performs 4--8 hop random walks over factual relations to construct multi-hop search prompts. These prompts are then solved using Tavily Search Engine (\url{https://www.tavily.com/}), with MiniMax 2.1 serving as the teacher model.

In addition, we include a new search dataset: a commercially cleared subset of OpenResearcher~\citep{li2026openresearcher}, a public SFT dataset designed for long-horizon research agents. OpenResearcher synthesizes over $97$K trajectories using \texttt{gpt-oss-120b}~\citep{agarwal2025gpt} as the teacher model in a fully offline browser environment. The environment is backed by a local search index over $15$M FineWeb documents~\citep{penedo2024fineweb} augmented with bootstrapped evidence documents, and exposes three structured browser tools: \texttt{search} for retrieving candidate pages, \texttt{open} for inspecting full document contents, and \texttt{find} for locating exact textual evidence within an opened document. The resulting trajectories capture long-horizon reasoning-action-observation loops in which the model iteratively decomposes research questions, gathers and inspects sources, localizes supporting evidence, and synthesizes grounded final answers. For \ourmodel, we did not regenerate the OpenResearcher data; instead, we curated the commercial-OK portion by removing examples whose source licensing was not cleared for commercial use. This filtering yields approximately $21.7$K SFT trajectories while preserving the original browser-tool interaction format for training open-ended research and evidence-grounded search behavior.

Finally, we work with data vendors to curate particularly challenging samples that require 50--100 searches, and collect SFT trajectories in our BrowseComp harness, described in Appendix~\ref{appendix:eval}. For these trajectories, we use MiniMax 2.5 and GLM 5.1 as teacher models.


\paragraph{Terminal-Use Capabilities.}
% \todo{Khushi / Charlie / Venkat}
To develop strong terminal-use capabilities, we constructed a large-scale dataset of synthetic agentic trajectories covering a broad range of terminal tasks, including software engineering, data processing, file operations, and scientific computing. Seed instructions were sourced from a combination of publicly available datasets: OpenCodeReasoning \citep{nvidia2025opencodeReasoning}, OpenMathReasoning \citep{nvidia2025openmathReasoning}, SWE-bench \citep{jimenez2024swebenchlanguagemodelsresolve}, SWE-Fixer-Train-110K \citep{swefixer2025}, SWE-rebench \citep{badertdinov2025swerebenchautomatedpipelinetask}, and SWE-smith \citep{pan2025swesmith}. Tasks were assembled from two complementary sources: one portion was adapted from existing math and coding SFT data in Nemotron-Cascade, reformatted to align with the terminal-use environment, while the other was synthetically generated using DeepSeek-V3.2 \citep{deepseekai2025deepseekv3technicalreport} to cover a wider variety of terminal scenarios not well-represented in existing benchmarks. For trajectory generation, DeepSeek-V3.2 was used as the acting agent within the Terminus-2 agent provided by the Harbor framework \citep{harbor2025}. Each trajectory unfolds over multiple episodes in which the model receives an initial task prompt and subsequently interacts with a live terminal environment, issuing commands and incorporating execution feedback to make iterative progress toward task completion. This multi-step, environment-grounded collection process encourages the model to learn practical tool use, adaptive planning, and error recovery behaviors that are characteristic of effective terminal operation. The final dataset comprises approximately 370K multi-turn conversations, consisting of a mixture of reasoning and non-reasoning trajectories.

\paragraph{Conversational Tool Use Capabilities.}
We scale conversational tool-use data via a fully
synthetic, six-stage generation pipeline. This includes User and Environment simulation. This pipeline is similar to what is described in Nemotron-3 Super. 

\paragraph{Software Issue Resolution.}
% \todo{@Sugam / Som / Wasi - SWE data}
To distill robust problem solving capabilities, we curated a diverse dataset of synthetic agent trajectories focused on resolving real-world GitHub issues. To this end, we generated synthetic trajectories by leveraging two distinct modeling paradigms: the reasoning-based (thinking) Minimax-M2.5 \citep{minimax2026m25} and the instruction-based (non-thinking) Qwen3-Coder-480B-A35B-Instruct \citep{yang2025qwen3technicalreport} for issue resolution. The underlying issue statements were curated from a diverse set of publicly available datasets, including SWE-Gym \citep{pan2025trainingsoftwareengineeringagents}, R2E-Gym \citep{jain2025r2egymproceduralenvironmentshybrid}, SWE-rebench \citep{badertdinov2025swerebenchautomatedpipelinetask}, and SWE-rebench-V2 \citep{badertdinov2026swe}. These trajectories were captured using the OpenHands \citep{wang2025openhandsopenplatformai}, SWE-agent \citep{yang2024sweagent}, and Mini-SWE-agent and Opencode harnesses. Raw agent rollouts contain many patterns that, while completing the task, would teach undesirable behaviors if used directly as SFT data, so we filter the rollout pool with a per-trajectory heuristic analyzer evaluates a fixed set of signals to make an include/exclude decision: submission
  integrity (trajectory must end with a valid submission action), disallowed git operations (any use of push, pull, fetch, clone, cherry-pick, reflog, fsck, remote, or ls-remote), edit-test loop anti-patterns (the model thrashes between edits and test runs without converging — either reverting its own edits or repeating the edit → test → edit cycle indefinitely), lost-in-exploration (the model spends most of its turns reading/searching the repo but rarely edits, indicating failure to localize and commit to a fix), tool-call hygiene (high rate of malformed tool calls), debug-artifact detection in the final patch (print(, pdb, breakpoint()), and a verification check that flags trajectories that edit but never run tests. This multi-model, multi-framework approach ensures the resulting dataset promotes generalization across both diverse problem-solving modes and varied agentic environments. 

\paragraph{Math / Proof Data.}
We use the Nemotron-Cascade-2~\citep{yang2026nemotron} math data, which sources problems primarily from Nemotron-Cascade~\citep{wang2025nemotroncascadescalingcascadedreinforcement,liu2025acereason} and Nemotron-Math-V2~\citep{du2025nemotronmath}.
For the non-proof math data, we collect 1.8M tool-calling samples and 1.9M non-tool samples, with responses generated by DeepSeek-V3.2 and DeepSeek-V3.2-Speciale, respectively.
For the mathematical natural language proof data, we source proof problems from the AOPS split of Nemotron-Math-Proofs-v1~\citep{du2025nemotronmath} and generate responses from DeepSeek-V3.2-Speciale for capabilities including proof generation, proof verification, and proof refinement.

\paragraph{Science.}
We prepare our science SFT data following the Nemotron Nano recipe~\citep{nvidia2025nemotron3nanoopen}, which combines synthetic, real-world, and document-derived seed data spanning physics, chemistry, and biology. Prompts and formats are diversified using NeMo Data Designer~\citep{nemo-data-designer}, and the resulting samples are filtered with an LLM judge for format compliance and reasoning quality. In addition to the Nano recipe, we generate web-search and web-search-with-Python reasoning traces using DeepSeek-V3.2~\citep{deepseekai2024deepseekv32}. For web-search traces, the model is provided access to the Tavily search engine (\url{https://www.tavily.com/}); for web-search-with-Python traces, the model additionally has access to a Python execution environment.

\paragraph{Chat.}
We create multi-turn chat SFT data by starting from seed prompts drawn from open conversational datasets such as LMArena \citep{chiang2024chatbot} and WildChat \citep{li2024wildchat}. For each prompt, we sample multiple candidate responses from GLM-5 \citep{glm5team2026glm5vibecodingagentic} and use Nemotron-GenRM \citep{Wang2025HelpSteer3Preference} to select the highest-quality response for the current turn. To extend the data into multi-turn conversations, we simulate the user with the same LLM under controlled prompting. The simulated user is guided by hand-crafted conversation strategies, such as building on prior content, asking for clarification, challenging assumptions, reframing the task, or applying the answer to a new context, in order to produce diverse and realistic dialogue trajectories. To build multi-turn robustness, we construct examples in which earlier assistant turns may contain suboptimal responses rather than the best-ranked candidates. The model is trained only on the final assistant response in SFT to produce high-quality current-turn outputs while remaining robust to imperfect previous turns.

\paragraph{Code.}
We collect our coding problems from modern competitive programming platforms such as Codeforces, AtCoder, AIZU, and CodeChef. Following \citet{yang2026nemotron}, we applied strict deduplication and aggressive filtering to improve the data quality and balance the problem difficulties. We choose GPT-OSS-120B~\citep{agarwal2025gpt} as our teacher model with the consideration of their strong reasoning ability and good verbosity, and apply the rejection sampling on those reasoning traces. This pipeline yields a final dataset comprising 1.2M Python reasoning traces, 1.0M C++14 reasoning traces, and 1.3M Python tool-calling reasoning traces for competitive coding.

\paragraph{CUDA.}
%@Jingquan Wang US - Cuda data
We constructed a large-scale synthetic CUDA dataset comprising approximately 100K samples for kernel generation, repair, and optimization. The dataset was built using an LLM-based synthetic data generation pipeline with DeepSeek-R1 and GPT-OSS-120B. Seed questions were sourced from popular open-source libraries, NVIDIA library API surfaces, and BackendBench~\citep{saroufimbackendbench}. The construction was motivated by the need for CUDA-specialized training data that reflects real-world GPU-kernel programming and benchmarking challenges, as highlighted by recent CUDA benchmark efforts such as SOL-ExecBench~\citep{lin2026solexecbenchspeedoflightbenchmarkingrealworld}. These seeds were used to generate two types of samples: \emph{PyTorch-reference-to-CUDA-kernel} samples and \emph{natural-language-specification-to-CUDA-kernel} samples, with each sample accompanied by reasoning. For each seed item, multiple candidate kernels were produced through LLM-based synthetic generation and rejection sampling. Candidate kernels were validated in an internal CUDA evaluation environment using compilation checks, numerical correctness tests, and runtime benchmarking. Invalid, non-compiling, or incorrect candidates were rejected. Among the remaining valid candidates, the best-performing kernel according to benchmarked runtime was retained. In addition, we collected traces from an internal CUDA agent to produce repair and optimization data. The repair samples contain a PyTorch reference, a faulty CUDA C++ kernel, the corresponding error message, and a corrected CUDA C++ kernel. The optimization samples contain a PyTorch reference, a slow CUDA C++ kernel, an Nsight Compute log, and an optimized CUDA C++ kernel. Beyond direct CUDA C++ kernel generation, we also generated CUDA-X library  data using publicly available documentation and official code samples. Following the same formulation as the CUDA-C data, we generated PyTorch references, aligned natural-language specifications, corresponding CUDA-X library implementations, and reasoning. The covered libraries include Thrust, CUB, cuBLAS, cuDNN, cuSPARSE, cuRAND, and cuSOLVER.

\paragraph{RTL.}
%@Chenhui Deng US RTL data
We use the ACE-RTL~\citep{deng2026ace} training data, which covers three major RTL task categories: specification-to-RTL generation, code editing, and code debugging. ACE-RTL builds on the seed RTL corpus from ScaleRTL~\citep{deng2025scalertl}, where seed designs are collected from license-checked open-source RTL repositories and processed through deduplication, filtering, and syntax validation before being used for data generation. For specification-to-RTL tasks, DeepSeek-R1 and GPT-OSS-120B are used to synthesize natural-language specifications paired with the corresponding golden RTL implementations derived from the seed designs. For editing and debugging tasks, the original seed RTL is treated as the golden implementation, while simplified variants with missing functionality or buggy variants with realistic injected design errors are generated as inputs; the corresponding specifications describe either the required feature extensions or diagnostic information needed to recover the golden RTL. After further filtering through syntax checks, benchmark decontamination, and semantic-alignment evaluation using human-defined rubrics, the final dataset contains around 1.2M training samples.

\paragraph{Multilingual.}
Our multilingual post-training data is a mixture of sentence-level parallel corpora and synthetic translations of English SFT examples in math, code, and science domains.
For the synthetic data, driven by data quality issues identified in the previous line-by-line translation pipeline, we introduced a new end-to-end translation pipeline that takes the full JSON object as input and generates the full JSON object as output.
Because this pipeline requires strong long-context capabilities from the translation model, we transitioned to \texttt{DeepSeek-V3-0324}.
After translation, we first perform heuristic format checking to ensure all outputs conform to the specified JSON format, followed by the same data filtration steps and light-weight post-editing step introduced in the Nemotron Super-V3 recipe~\citep{nvidia2026nemotron3superopen}.
Ablation studies on Japanese data showed significant improvements measured by MMLU-ProX~\citep{xuan2025mmluprox}, confirming quality improvements over the line-by-line pipeline.
For Nemotron Ultra-V3, we synthesized data for Hindi, Japanese, Korean, and Brazilian Portuguese using this pipeline, while reusing existing multilingual synthetic data from Super-V3 for the remaining languages.

\subsubsection{Data Packing}
% @Wenliang - Blending / packing details
To train efficiently on a large and diverse collection of datasets, we adopt a length-aware best-fit packing strategy~\citep{ding2024fewer}, which packs multiple conversations into sequences up to a maximum context length. 
Our packer is highly memory-efficient: it reads and interleaves all source files in a round-robin fashion, maintains only a fixed-size pool of open sequences in memory, and retires a sequence once its residual capacity falls below a small tolerance. 
Each incoming conversation is assigned to the partially filled sequence whose remaining capacity it most tightly fits, following a best-fit rule that minimizes padding overhead. 
We neither truncate nor split conversations, preserving complete context and reducing hallucinations. 
Additionally, we enforce an in-pack deduplication constraint to prevent identical prompts from co-occurring within the same sequence. 
After packing, we perform a final shuffle over all completed packs. 
Overall, our implementation ensures that each packed sequence draws from a broad, well-mixed cross-section of the data distribution rather than clustering examples from any single source. 
This thorough mixing is essential at scale, where naive concatenation of per-source shards would otherwise induce strong distributional locality and degrade optimization stability.